% Options for packages loaded elsewhere
\PassOptionsToPackage{unicode}{hyperref}
\PassOptionsToPackage{hyphens}{url}
%
\documentclass[
]{article}
\usepackage{amsmath,amssymb}
\usepackage{iftex}
\ifPDFTeX
  \usepackage[T1]{fontenc}
  \usepackage[utf8]{inputenc}
  \usepackage{textcomp} % provide euro and other symbols
\else % if luatex or xetex
  \usepackage{unicode-math} % this also loads fontspec
  \defaultfontfeatures{Scale=MatchLowercase}
  \defaultfontfeatures[\rmfamily]{Ligatures=TeX,Scale=1}
\fi
\usepackage{lmodern}
\ifPDFTeX\else
  % xetex/luatex font selection
\fi
% Use upquote if available, for straight quotes in verbatim environments
\IfFileExists{upquote.sty}{\usepackage{upquote}}{}
\IfFileExists{microtype.sty}{% use microtype if available
  \usepackage[]{microtype}
  \UseMicrotypeSet[protrusion]{basicmath} % disable protrusion for tt fonts
}{}
\makeatletter
\@ifundefined{KOMAClassName}{% if non-KOMA class
  \IfFileExists{parskip.sty}{%
    \usepackage{parskip}
  }{% else
    \setlength{\parindent}{0pt}
    \setlength{\parskip}{6pt plus 2pt minus 1pt}}
}{% if KOMA class
  \KOMAoptions{parskip=half}}
\makeatother
\usepackage{xcolor}
\usepackage[margin=1in]{geometry}
\usepackage{color}
\usepackage{fancyvrb}
\newcommand{\VerbBar}{|}
\newcommand{\VERB}{\Verb[commandchars=\\\{\}]}
\DefineVerbatimEnvironment{Highlighting}{Verbatim}{commandchars=\\\{\}}
% Add ',fontsize=\small' for more characters per line
\usepackage{framed}
\definecolor{shadecolor}{RGB}{248,248,248}
\newenvironment{Shaded}{\begin{snugshade}}{\end{snugshade}}
\newcommand{\AlertTok}[1]{\textcolor[rgb]{0.94,0.16,0.16}{#1}}
\newcommand{\AnnotationTok}[1]{\textcolor[rgb]{0.56,0.35,0.01}{\textbf{\textit{#1}}}}
\newcommand{\AttributeTok}[1]{\textcolor[rgb]{0.13,0.29,0.53}{#1}}
\newcommand{\BaseNTok}[1]{\textcolor[rgb]{0.00,0.00,0.81}{#1}}
\newcommand{\BuiltInTok}[1]{#1}
\newcommand{\CharTok}[1]{\textcolor[rgb]{0.31,0.60,0.02}{#1}}
\newcommand{\CommentTok}[1]{\textcolor[rgb]{0.56,0.35,0.01}{\textit{#1}}}
\newcommand{\CommentVarTok}[1]{\textcolor[rgb]{0.56,0.35,0.01}{\textbf{\textit{#1}}}}
\newcommand{\ConstantTok}[1]{\textcolor[rgb]{0.56,0.35,0.01}{#1}}
\newcommand{\ControlFlowTok}[1]{\textcolor[rgb]{0.13,0.29,0.53}{\textbf{#1}}}
\newcommand{\DataTypeTok}[1]{\textcolor[rgb]{0.13,0.29,0.53}{#1}}
\newcommand{\DecValTok}[1]{\textcolor[rgb]{0.00,0.00,0.81}{#1}}
\newcommand{\DocumentationTok}[1]{\textcolor[rgb]{0.56,0.35,0.01}{\textbf{\textit{#1}}}}
\newcommand{\ErrorTok}[1]{\textcolor[rgb]{0.64,0.00,0.00}{\textbf{#1}}}
\newcommand{\ExtensionTok}[1]{#1}
\newcommand{\FloatTok}[1]{\textcolor[rgb]{0.00,0.00,0.81}{#1}}
\newcommand{\FunctionTok}[1]{\textcolor[rgb]{0.13,0.29,0.53}{\textbf{#1}}}
\newcommand{\ImportTok}[1]{#1}
\newcommand{\InformationTok}[1]{\textcolor[rgb]{0.56,0.35,0.01}{\textbf{\textit{#1}}}}
\newcommand{\KeywordTok}[1]{\textcolor[rgb]{0.13,0.29,0.53}{\textbf{#1}}}
\newcommand{\NormalTok}[1]{#1}
\newcommand{\OperatorTok}[1]{\textcolor[rgb]{0.81,0.36,0.00}{\textbf{#1}}}
\newcommand{\OtherTok}[1]{\textcolor[rgb]{0.56,0.35,0.01}{#1}}
\newcommand{\PreprocessorTok}[1]{\textcolor[rgb]{0.56,0.35,0.01}{\textit{#1}}}
\newcommand{\RegionMarkerTok}[1]{#1}
\newcommand{\SpecialCharTok}[1]{\textcolor[rgb]{0.81,0.36,0.00}{\textbf{#1}}}
\newcommand{\SpecialStringTok}[1]{\textcolor[rgb]{0.31,0.60,0.02}{#1}}
\newcommand{\StringTok}[1]{\textcolor[rgb]{0.31,0.60,0.02}{#1}}
\newcommand{\VariableTok}[1]{\textcolor[rgb]{0.00,0.00,0.00}{#1}}
\newcommand{\VerbatimStringTok}[1]{\textcolor[rgb]{0.31,0.60,0.02}{#1}}
\newcommand{\WarningTok}[1]{\textcolor[rgb]{0.56,0.35,0.01}{\textbf{\textit{#1}}}}
\usepackage{graphicx}
\makeatletter
\newsavebox\pandoc@box
\newcommand*\pandocbounded[1]{% scales image to fit in text height/width
  \sbox\pandoc@box{#1}%
  \Gscale@div\@tempa{\textheight}{\dimexpr\ht\pandoc@box+\dp\pandoc@box\relax}%
  \Gscale@div\@tempb{\linewidth}{\wd\pandoc@box}%
  \ifdim\@tempb\p@<\@tempa\p@\let\@tempa\@tempb\fi% select the smaller of both
  \ifdim\@tempa\p@<\p@\scalebox{\@tempa}{\usebox\pandoc@box}%
  \else\usebox{\pandoc@box}%
  \fi%
}
% Set default figure placement to htbp
\def\fps@figure{htbp}
\makeatother
\setlength{\emergencystretch}{3em} % prevent overfull lines
\providecommand{\tightlist}{%
  \setlength{\itemsep}{0pt}\setlength{\parskip}{0pt}}
\setcounter{secnumdepth}{-\maxdimen} % remove section numbering
% definitions for citeproc citations
\NewDocumentCommand\citeproctext{}{}
\NewDocumentCommand\citeproc{mm}{%
  \begingroup\def\citeproctext{#2}\cite{#1}\endgroup}
\makeatletter
 % allow citations to break across lines
 \let\@cite@ofmt\@firstofone
 % avoid brackets around text for \cite:
 \def\@biblabel#1{}
 \def\@cite#1#2{{#1\if@tempswa , #2\fi}}
\makeatother
\newlength{\cslhangindent}
\setlength{\cslhangindent}{1.5em}
\newlength{\csllabelwidth}
\setlength{\csllabelwidth}{3em}
\newenvironment{CSLReferences}[2] % #1 hanging-indent, #2 entry-spacing
 {\begin{list}{}{%
  \setlength{\itemindent}{0pt}
  \setlength{\leftmargin}{0pt}
  \setlength{\parsep}{0pt}
  % turn on hanging indent if param 1 is 1
  \ifodd #1
   \setlength{\leftmargin}{\cslhangindent}
   \setlength{\itemindent}{-1\cslhangindent}
  \fi
  % set entry spacing
  \setlength{\itemsep}{#2\baselineskip}}}
 {\end{list}}
\usepackage{calc}
\newcommand{\CSLBlock}[1]{\hfill\break\parbox[t]{\linewidth}{\strut\ignorespaces#1\strut}}
\newcommand{\CSLLeftMargin}[1]{\parbox[t]{\csllabelwidth}{\strut#1\strut}}
\newcommand{\CSLRightInline}[1]{\parbox[t]{\linewidth - \csllabelwidth}{\strut#1\strut}}
\newcommand{\CSLIndent}[1]{\hspace{\cslhangindent}#1}
\usepackage{booktabs}
\usepackage{caption}
\usepackage{longtable}
\usepackage{colortbl}
\usepackage{array}
\usepackage{anyfontsize}
\usepackage{multirow}
\usepackage{bookmark}
\IfFileExists{xurl.sty}{\usepackage{xurl}}{} % add URL line breaks if available
\urlstyle{same}
\hypersetup{
  pdftitle={Análise dos Fatores Associados à Resposta Terapêutica Oncológica},
  pdfauthor={Márcio Souza},
  hidelinks,
  pdfcreator={LaTeX via pandoc}}

\title{Análise dos Fatores Associados à Resposta Terapêutica Oncológica}
\author{Márcio Souza}
\date{2025-11-22}

\begin{document}
\maketitle

\subsection{Preparando os dados}\label{preparando-os-dados}

\begin{Shaded}
\begin{Highlighting}[]
\NormalTok{dados }\OtherTok{\textless{}{-}} \FunctionTok{read\_xls}\NormalTok{(}\StringTok{"base\_final\_final\_com\_dicionário\_das\_variáveis.xls"}\NormalTok{)}
\FunctionTok{names}\NormalTok{(dados)[}\DecValTok{29}\NormalTok{] }\OtherTok{=} \StringTok{"descricao"}

\CommentTok{\# Criando o data.frame com a tabela de estadiamento}
\NormalTok{tabela\_estadiamento }\OtherTok{\textless{}{-}} \FunctionTok{data.frame}\NormalTok{(}
  \AttributeTok{estadiam =} \FunctionTok{c}\NormalTok{(}\StringTok{"0"}\NormalTok{, }\StringTok{"1"}\NormalTok{, }\StringTok{"1C"}\NormalTok{, }\StringTok{"2"}\NormalTok{, }\StringTok{"2A"}\NormalTok{, }\StringTok{"2B"}\NormalTok{, }\StringTok{"2C"}\NormalTok{, }
             \StringTok{"3"}\NormalTok{, }\StringTok{"3A"}\NormalTok{, }\StringTok{"3B"}\NormalTok{, }\StringTok{"3C"}\NormalTok{, }\StringTok{"4"}\NormalTok{, }\StringTok{"4A"}\NormalTok{, }\StringTok{"4B"}\NormalTok{, }
             \StringTok{"88"}\NormalTok{, }\StringTok{"99"}\NormalTok{),}
  
  \AttributeTok{rotulo =} \FunctionTok{c}\NormalTok{(}\StringTok{"Estádio 0"}\NormalTok{, }\StringTok{"Estádio I"}\NormalTok{, }\StringTok{"Estádio IC"}\NormalTok{, }
             \StringTok{"Estádio II"}\NormalTok{, }\StringTok{"Estádio IIA"}\NormalTok{, }\StringTok{"Estádio IIB"}\NormalTok{, }\StringTok{"Estádio IIC"}\NormalTok{, }
             \StringTok{"Estádio III"}\NormalTok{, }\StringTok{"Estádio IIIA"}\NormalTok{, }\StringTok{"Estádio IIIB"}\NormalTok{, }\StringTok{"Estádio IIIC"}\NormalTok{, }
             \StringTok{"Estádio IV"}\NormalTok{, }\StringTok{"Estádio IVA"}\NormalTok{, }\StringTok{"Estádio IVB"}\NormalTok{, }
             \StringTok{"Não se aplica"}\NormalTok{, }\StringTok{"Ignorado"}\NormalTok{),}
  
  \CommentTok{\# {-}{-}{-} NOVA COLUNA 1: Agrupamento Clássico (Romano) {-}{-}{-}}
  \AttributeTok{estadio\_geral =} \FunctionTok{c}\NormalTok{(}\StringTok{"0"}\NormalTok{, }\StringTok{"I"}\NormalTok{, }\StringTok{"I"}\NormalTok{, }
                    \StringTok{"II"}\NormalTok{, }\StringTok{"II"}\NormalTok{, }\StringTok{"II"}\NormalTok{, }\StringTok{"II"}\NormalTok{, }
                    \StringTok{"III"}\NormalTok{, }\StringTok{"III"}\NormalTok{, }\StringTok{"III"}\NormalTok{, }\StringTok{"III"}\NormalTok{, }
                    \StringTok{"IV"}\NormalTok{, }\StringTok{"IV"}\NormalTok{, }\StringTok{"IV"}\NormalTok{, }
                    \StringTok{"N/A"}\NormalTok{, }\StringTok{"Ignorado"}\NormalTok{),}
  
  \CommentTok{\# {-}{-}{-} NOVA COLUNA 2: Agrupamento Epidemiológico (Binário) {-}{-}{-}}
  \CommentTok{\# Geralmente: 0, I e II são "Precoce/Inicial", III e IV são "Avançado"}
  \AttributeTok{estadiam\_bin =} \FunctionTok{c}\NormalTok{(}\StringTok{"Inicial"}\NormalTok{, }\StringTok{"Inicial"}\NormalTok{, }\StringTok{"Inicial"}\NormalTok{, }
                            \StringTok{"Inicial"}\NormalTok{, }\StringTok{"Inicial"}\NormalTok{, }\StringTok{"Inicial"}\NormalTok{, }\StringTok{"Inicial"}\NormalTok{, }
                            \StringTok{"Avançado"}\NormalTok{, }\StringTok{"Avançado"}\NormalTok{, }\StringTok{"Avançado"}\NormalTok{, }\StringTok{"Avançado"}\NormalTok{, }
                            \StringTok{"Avançado"}\NormalTok{, }\StringTok{"Avançado"}\NormalTok{, }\StringTok{"Avançado"}\NormalTok{, }
                            \ConstantTok{NA}\NormalTok{, }\ConstantTok{NA}\NormalTok{), }\CommentTok{\# 88 e 99 viram NA para não sujar a análise}
  
  \AttributeTok{descricao =} \FunctionTok{c}\NormalTok{(}
    \StringTok{"Carcinoma in situ (não invasivo)"}\NormalTok{,}
    \StringTok{"Tumor inicial, pequeno e localizado"}\NormalTok{,}
    \StringTok{"Subdivisão do Estádio I"}\NormalTok{,}
    \StringTok{"Tumor localmente avançado"}\NormalTok{,}
    \StringTok{"Subdivisão inicial do Estádio II"}\NormalTok{,}
    \StringTok{"Subdivisão intermediária do Estádio II"}\NormalTok{,}
    \StringTok{"Subdivisão avançada do Estádio II"}\NormalTok{,}
    \StringTok{"Tumor avançado regionalmente"}\NormalTok{,}
    \StringTok{"Subdivisão inicial do Estádio III"}\NormalTok{,}
    \StringTok{"Subdivisão intermediária do Estádio III"}\NormalTok{,}
    \StringTok{"Subdivisão avançada do Estádio III"}\NormalTok{,}
    \StringTok{"Metástase à distância"}\NormalTok{,}
    \StringTok{"Metástase ou avanço local agressivo"}\NormalTok{,}
    \StringTok{"Metástase à distância avançada"}\NormalTok{,}
    \StringTok{"Tumores sem estadiamento TNM"}\NormalTok{,}
    \StringTok{"Informação desconhecida ou não registrada"}
\NormalTok{  ),}
  
  \AttributeTok{stringsAsFactors =} \ConstantTok{FALSE}
\NormalTok{)}

\CommentTok{\# Converter variáveis de data}
\NormalTok{dados\_clean }\OtherTok{\textless{}{-}}\NormalTok{ dados }\SpecialCharTok{\%\textgreater{}\%}
  \FunctionTok{left\_join}\NormalTok{(tabela\_estadiamento[,}\FunctionTok{c}\NormalTok{(}\DecValTok{1}\NormalTok{,}\DecValTok{4}\NormalTok{)], }\StringTok{"estadiam"}\NormalTok{) }\SpecialCharTok{\%\textgreater{}\%} 
  \FunctionTok{filter}\NormalTok{(tratamentoantesdodiag }\SpecialCharTok{!=} \StringTok{"tratamentoantesdodiag"}\NormalTok{) }\SpecialCharTok{\%\textgreater{}\%}
\NormalTok{  dplyr}\SpecialCharTok{::}\FunctionTok{select}\NormalTok{(}\SpecialCharTok{{-}}\NormalTok{tratamentoantesdodiag,}\SpecialCharTok{{-}}\NormalTok{tpcaso,}\SpecialCharTok{{-}}\NormalTok{datainitrt,}\SpecialCharTok{{-}}\NormalTok{datapricon,}
         \SpecialCharTok{{-}}\NormalTok{trancursodiagtrat,}\SpecialCharTok{{-}}\NormalTok{trancurso1consultdiag,}\SpecialCharTok{{-}}\NormalTok{diagnosticoantesdaconsulta,}
         \SpecialCharTok{{-}}\NormalTok{dtdiagno,}\SpecialCharTok{{-}}\NormalTok{ocupacao,}\SpecialCharTok{{-}}\NormalTok{codigo,}\SpecialCharTok{{-}}\NormalTok{faixaconsultadiag,}\SpecialCharTok{{-}}\NormalTok{estadiam) }\SpecialCharTok{\%\textgreater{}\%} 
  \FunctionTok{mutate}\NormalTok{(}
    \AttributeTok{faixa\_etaria\_prim\_cons =} \FunctionTok{cut}\NormalTok{(}
\NormalTok{      idade1consulta,}
      \AttributeTok{breaks =} \FunctionTok{c}\NormalTok{(}\FunctionTok{seq}\NormalTok{(}\DecValTok{0}\NormalTok{, }\DecValTok{80}\NormalTok{, }\AttributeTok{by =} \DecValTok{10}\NormalTok{), }\ConstantTok{Inf}\NormalTok{), }\CommentTok{\# Define as quebras de 5 em 5}
      \AttributeTok{labels =} \FunctionTok{c}\NormalTok{(}
        \StringTok{"0{-}9"}\NormalTok{, }\StringTok{"10{-}19"}\NormalTok{, }\StringTok{"20{-}29"}\NormalTok{, }\StringTok{"30{-}39"}\NormalTok{,}
        \StringTok{"40{-}49"}\NormalTok{, }\StringTok{"50{-}59"}\NormalTok{, }\StringTok{"60{-}69"}\NormalTok{, }\StringTok{"70{-}79"}\NormalTok{,}
        \StringTok{"80+"}
\NormalTok{      ),}
      \AttributeTok{right =} \ConstantTok{FALSE} \CommentTok{\# Importante: Faz o intervalo ser [0, 5) {-}\textgreater{} inclui 0, exclui 5}
\NormalTok{    ),}
    \AttributeTok{origemenSUS =} \FunctionTok{ifelse}\NormalTok{(origemencaminhamento }\SpecialCharTok{==} \StringTok{"SUS"}\NormalTok{, }\StringTok{"Sim"}\NormalTok{, }\StringTok{"Não"}\NormalTok{),}
    \AttributeTok{raca =} \FunctionTok{case\_when}\NormalTok{(}
\NormalTok{      racacor2 }\SpecialCharTok{==} \StringTok{"branca"} \SpecialCharTok{\textasciitilde{}} \StringTok{"Branca"}\NormalTok{,}
\NormalTok{      racacor2 }\SpecialCharTok{==} \StringTok{"parda"} \SpecialCharTok{\textasciitilde{}} \StringTok{"Parda"}\NormalTok{,}
\NormalTok{      racacor2 }\SpecialCharTok{==} \StringTok{"preta"} \SpecialCharTok{\textasciitilde{}} \StringTok{"Preta"}\NormalTok{,}
      \ConstantTok{TRUE} \SpecialCharTok{\textasciitilde{}} \ConstantTok{NA\_character\_}
\NormalTok{    ),}
    \AttributeTok{basediag\_agrupada =} \FunctionTok{case\_when}\NormalTok{(}
      \CommentTok{\# Grupo Microscópico (Certeza)}
\NormalTok{      basediag }\SpecialCharTok{\%in\%} \FunctionTok{c}\NormalTok{(}\StringTok{"histologia tum prim"}\NormalTok{, }
                      \StringTok{"histologia metastase"}\NormalTok{, }
                      \StringTok{"citologia"}\NormalTok{) }\SpecialCharTok{\textasciitilde{}} \StringTok{"Confirmacao Microscopica"}\NormalTok{,}
      
      \CommentTok{\# Grupo Não Microscópico (Evidência Clínica/Imagem)}
\NormalTok{      basediag }\SpecialCharTok{\%in\%} \FunctionTok{c}\NormalTok{(}\StringTok{"clinica"}\NormalTok{, }
                      \StringTok{"exame por imagem"}\NormalTok{, }
                      \StringTok{"marcadores tumorais"}\NormalTok{, }
                      \StringTok{"pesquisa clinica"}\NormalTok{) }\SpecialCharTok{\textasciitilde{}} \StringTok{"Diagnostico Clinico/Imagem"}\NormalTok{,}
      
      \CommentTok{\# O que fazer com s/info? (Sugestão: tratar como NA ou Outros)}
\NormalTok{      basediag }\SpecialCharTok{==} \StringTok{"s/info"} \SpecialCharTok{\textasciitilde{}} \ConstantTok{NA\_character\_}\NormalTok{, }\CommentTok{\# Ou "Ignorado"}
      
      \ConstantTok{TRUE} \SpecialCharTok{\textasciitilde{}} \ConstantTok{NA\_character\_} \CommentTok{\# Segurança para casos não mapeados}
\NormalTok{    ),}
    \AttributeTok{status\_conjugal\_bin =} \FunctionTok{case\_when}\NormalTok{(}
\NormalTok{      estadocivil }\SpecialCharTok{\%in\%} \FunctionTok{c}\NormalTok{(}\StringTok{"casado"}\NormalTok{, }\StringTok{"uniao consensual"}\NormalTok{) }\SpecialCharTok{\textasciitilde{}} \StringTok{"Com Companheiro"}\NormalTok{,}
\NormalTok{      estadocivil }\SpecialCharTok{\%in\%} \FunctionTok{c}\NormalTok{(}\StringTok{"solteiro"}\NormalTok{, }\StringTok{"separacao judicial"}\NormalTok{, }\StringTok{"viuvo"}\NormalTok{) }\SpecialCharTok{\textasciitilde{}} \StringTok{"Sem Companheiro"}\NormalTok{,}
      \ConstantTok{TRUE} \SpecialCharTok{\textasciitilde{}} \ConstantTok{NA\_character\_}
\NormalTok{    ),}
    \AttributeTok{resposta\_tratamento =} \FunctionTok{case\_when}\NormalTok{(}
\NormalTok{      estadofinal }\SpecialCharTok{\%in\%} \FunctionTok{c}\NormalTok{(}\StringTok{"remissao completa"}\NormalTok{, }\StringTok{"remissao parcial"}\NormalTok{) }\SpecialCharTok{\textasciitilde{}} \StringTok{"Satisfatória"}\NormalTok{,}
      
\NormalTok{      estadofinal }\SpecialCharTok{\%in\%} \FunctionTok{c}\NormalTok{(}\StringTok{"doenca estavel"}\NormalTok{, }\StringTok{"doenca progressao"}\NormalTok{, }
                         \StringTok{"obito"}\NormalTok{, }\StringTok{"suporte terapeutico"}\NormalTok{) }\SpecialCharTok{\textasciitilde{}} \StringTok{"Insatisfatória"}\NormalTok{,}
      
      \CommentTok{\# Todo o resto vira NA}
      \ConstantTok{TRUE} \SpecialCharTok{\textasciitilde{}} \ConstantTok{NA\_character\_} 
\NormalTok{    ),}
    \AttributeTok{tipo\_exame\_agrupado =} \FunctionTok{case\_when}\NormalTok{(}
      \CommentTok{\# Grupo 1: Exame Bioquímico}
\NormalTok{      examesdiagn }\SpecialCharTok{==} \StringTok{"marcadores tumorais"} \SpecialCharTok{\textasciitilde{}} \StringTok{"Marcadores Tumorais"}\NormalTok{,}
      
      \CommentTok{\# Grupo 2: Exame Tecidual (Biópsia)}
\NormalTok{      examesdiagn }\SpecialCharTok{==} \StringTok{"anatomia patologica"} \SpecialCharTok{\textasciitilde{}} \StringTok{"Anatomia Patologica"}\NormalTok{,}
      
      \CommentTok{\# Grupo 3: Todo o resto vira NA (ou "Outros" se preferir não perder linhas)}
      \CommentTok{\# Optei por NA pois \textless{}100 casos em 24mil não geram inferência segura}
      \ConstantTok{TRUE} \SpecialCharTok{\textasciitilde{}} \ConstantTok{NA\_character\_} 
\NormalTok{    ),}
    \AttributeTok{escolaridade\_agrupada =} \FunctionTok{case\_when}\NormalTok{(}
      \CommentTok{\# Nível 1: Baixa (Analfabeto funcional ou Primário incompleto)}
\NormalTok{      escolaridade }\SpecialCharTok{\%in\%} \FunctionTok{c}\NormalTok{(}\StringTok{"nenhuma"}\NormalTok{, }\StringTok{"fund. incompleto"}\NormalTok{) }\SpecialCharTok{\textasciitilde{}} \StringTok{"Baixa Escolaridade"}\NormalTok{,}
      
      \CommentTok{\# Nível 2: Média (Fundamental completo até Médio)}
\NormalTok{      escolaridade }\SpecialCharTok{\%in\%} \FunctionTok{c}\NormalTok{(}\StringTok{"fund. completo"}\NormalTok{, }\StringTok{"nivel medio"}\NormalTok{) }\SpecialCharTok{\textasciitilde{}} \StringTok{"Media Escolaridade"}\NormalTok{,}
      
      \CommentTok{\# Nível 3: Alta (Entrou na faculdade)}
\NormalTok{      escolaridade }\SpecialCharTok{\%in\%} \FunctionTok{c}\NormalTok{(}\StringTok{"superior incompleto"}\NormalTok{, }\StringTok{"superior completo"}\NormalTok{) }\SpecialCharTok{\textasciitilde{}} \StringTok{"Alta Escolaridade"}\NormalTok{,}
      
      \CommentTok{\# Tratamento de Missing}
\NormalTok{      escolaridade }\SpecialCharTok{==} \StringTok{"s/info"} \SpecialCharTok{\textasciitilde{}} \ConstantTok{NA\_character\_}\NormalTok{,}
      \ConstantTok{TRUE} \SpecialCharTok{\textasciitilde{}} \ConstantTok{NA\_character\_}
\NormalTok{    ),}
    \AttributeTok{tratamento\_agrupado =} \FunctionTok{case\_when}\NormalTok{(}
      \CommentTok{\# Os três grandes pilares do tratamento}
\NormalTok{      primeirotrathosp }\SpecialCharTok{==} \StringTok{"cirurgia"} \SpecialCharTok{\textasciitilde{}} \StringTok{"Cirurgia"}\NormalTok{,}
\NormalTok{      primeirotrathosp }\SpecialCharTok{==} \StringTok{"radioterapia"} \SpecialCharTok{\textasciitilde{}} \StringTok{"Radioterapia"}\NormalTok{,}
\NormalTok{      primeirotrathosp }\SpecialCharTok{==} \StringTok{"hormonioterapia"} \SpecialCharTok{\textasciitilde{}} \StringTok{"Hormonioterapia"}\NormalTok{,}
      
      \CommentTok{\# O grande grupo dos "Não tratados" (Vigilância ou Encaminhados)}
\NormalTok{      primeirotrathosp }\SpecialCharTok{\%in\%} \FunctionTok{c}\NormalTok{(}\StringTok{"NTI"}\NormalTok{, }\StringTok{"nenhum"}\NormalTok{) }\SpecialCharTok{\textasciitilde{}} \StringTok{"Sem Tratamento Inicial"}\NormalTok{,}
      
      \CommentTok{\# Agrupando o "Resto" (Quimio é minoria aqui, Imuno é traço)}
\NormalTok{      primeirotrathosp }\SpecialCharTok{\%in\%} \FunctionTok{c}\NormalTok{(}\StringTok{"quimioterapia"}\NormalTok{, }\StringTok{"imunoterapia"}\NormalTok{, }\StringTok{"outros"}\NormalTok{) }\SpecialCharTok{\textasciitilde{}} \StringTok{"Quimio/Outros"}\NormalTok{,}
      
      \CommentTok{\# Tratamento de Missing}
\NormalTok{      primeirotrathosp }\SpecialCharTok{==} \StringTok{"s/info"} \SpecialCharTok{\textasciitilde{}} \ConstantTok{NA\_character\_}\NormalTok{,}
      \ConstantTok{TRUE} \SpecialCharTok{\textasciitilde{}} \ConstantTok{NA\_character\_}
\NormalTok{    ),}
    \CommentTok{\# Normalizar texto para minusculo e sem acentos ajuda no \textquotesingle{}match\textquotesingle{}}
    \AttributeTok{ocupacao\_temp =} \FunctionTok{iconv}\NormalTok{(}\FunctionTok{tolower}\NormalTok{(descricao), }\AttributeTok{to=}\StringTok{"ASCII//TRANSLIT"}\NormalTok{),}
    
    \AttributeTok{classe\_trabalho =} \FunctionTok{case\_when}\NormalTok{(}
      
      \CommentTok{\# {-}{-}{-} GRUPO 1: RURAL / AGRO (Prioridade alta para nao confundir com manual urbano) {-}{-}{-}}
      \FunctionTok{str\_detect}\NormalTok{(ocupacao\_temp, }\StringTok{"agro|rural|lavrador|pesca|agricola|florestal|cultura de|pecuaria|trabalhador da cultura"}\NormalTok{) }\SpecialCharTok{\textasciitilde{}} \StringTok{"Rural/Agropecuaria"}\NormalTok{,}
      
      \CommentTok{\# {-}{-}{-} GRUPO 2: ELITE / SUPERIOR / DIRIGENTES {-}{-}{-}}
      \FunctionTok{str\_detect}\NormalTok{(ocupacao\_temp, }\StringTok{"engenheir|medic|advogad|diretor|gerente|professo|arquitet|economista|analista|biolog|farmaceut|dentista|oficiais|universit|geolog|superiores|empresario|pesquisador|quimic"}\NormalTok{) }\SpecialCharTok{\textasciitilde{}} \StringTok{"Superior/Dirigente"}\NormalTok{,}
      
      \CommentTok{\# {-}{-}{-} GRUPO 3: MANUAL / OPERACIONAL / INDUSTRIAL {-}{-}{-}}
      \FunctionTok{str\_detect}\NormalTok{(ocupacao\_temp, }\StringTok{"pedreiro|servente|braçal|operador|mecanic|metalurg|soldador|montador|eletricista|pintor|carpinteiro|marceneiro|costureir|maquinista|motorista|condutor|ajustador|ferramenteiro|tecelao|trefilador|vidreiro|sapateiro|padeiro|acougueiro"}\NormalTok{) }\SpecialCharTok{\textasciitilde{}} \StringTok{"Manual/Operacional"}\NormalTok{,}
      
      \CommentTok{\# {-}{-}{-} GRUPO 4: TÉCNICOS E ADMINISTRATIVOS {-}{-}{-}}
      \FunctionTok{str\_detect}\NormalTok{(ocupacao\_temp, }\StringTok{"tecnico|administra|escritorio|secretari|auxiliar|contabil|bancario|recepcionista|desenhista|fiscal|inspetor|agente|assistente|datilografo"}\NormalTok{) }\SpecialCharTok{\textasciitilde{}} \StringTok{"Tecnico/Administrativo"}\NormalTok{,}
      
      \CommentTok{\# {-}{-}{-} GRUPO 5: SERVIÇOS / COMÉRCIO / SEGURANÇA {-}{-}{-}}
      \FunctionTok{str\_detect}\NormalTok{(ocupacao\_temp, }\StringTok{"vendedor|comerci|atendente|vigia|policia|bombeiro|militar|seguranca|zelador|porteiro|limpeza|domestica|cozinheir|garcom|barbeiro|cabeleireiro|manicure|frentista|carteiro|telefonista"}\NormalTok{) }\SpecialCharTok{\textasciitilde{}} \StringTok{"Servicos/Comercio"}\NormalTok{,}
      
      \CommentTok{\# {-}{-}{-} TRATAMENTO DE MISSING / RESTO {-}{-}{-}}
      \FunctionTok{str\_detect}\NormalTok{(ocupacao\_temp, }\StringTok{"nao se aplica|sem info|ignorado"}\NormalTok{) }\SpecialCharTok{\textasciitilde{}} \ConstantTok{NA\_character\_}\NormalTok{,}
      
      \CommentTok{\# Categoria de segurança para o que sobrar}
      \ConstantTok{TRUE} \SpecialCharTok{\textasciitilde{}} \StringTok{"Outros/Nao Classificado"}
\NormalTok{    ),}
    \AttributeTok{motivo\_nao\_tratamento =} \FunctionTok{case\_when}\NormalTok{(}
\NormalTok{      rzntr2 }\SpecialCharTok{\%in\%} \FunctionTok{c}\NormalTok{(}\StringTok{"abandono"}\NormalTok{, }\StringTok{"recusa"}\NormalTok{) }\SpecialCharTok{\textasciitilde{}} \StringTok{"Recusa/Abandono"}\NormalTok{,}
\NormalTok{      rzntr2 }\SpecialCharTok{\%in\%} \FunctionTok{c}\NormalTok{(}\StringTok{"doenca avancada"}\NormalTok{, }\StringTok{"obito"}\NormalTok{, }\StringTok{"complicacoes"}\NormalTok{) }\SpecialCharTok{\textasciitilde{}} \StringTok{"Condicao Clinica Ruim"}\NormalTok{,}
\NormalTok{      rzntr2 }\SpecialCharTok{==} \StringTok{"realizado fora"} \SpecialCharTok{\textasciitilde{}} \StringTok{"Transferido"}\NormalTok{,}
      \ConstantTok{TRUE} \SpecialCharTok{\textasciitilde{}} \StringTok{"Tratado/Nao se aplica"}
\NormalTok{    ),}
    \AttributeTok{alcoolismo\_agrupado =} \FunctionTok{case\_when}\NormalTok{(}
      \CommentTok{\# Grupo de Referência (Protegido/Baixo Risco)}
\NormalTok{      alcoolismo }\SpecialCharTok{==} \StringTok{"nunca"} \SpecialCharTok{\textasciitilde{}} \StringTok{"Nunca Bebeu"}\NormalTok{,}
      
      \CommentTok{\# Grupo de Risco Passado}
\NormalTok{      alcoolismo }\SpecialCharTok{==} \StringTok{"ex{-}consumidor"} \SpecialCharTok{\textasciitilde{}} \StringTok{"Ex{-}Consumidor"}\NormalTok{,}
      
      \CommentTok{\# Grupo de Risco Ativo}
\NormalTok{      alcoolismo }\SpecialCharTok{==} \StringTok{"sim"} \SpecialCharTok{\textasciitilde{}} \StringTok{"Consumidor Ativo"}\NormalTok{,}
      
      \CommentTok{\# Tratamento de Missing (s/info, nao avaliado, nti)}
      \ConstantTok{TRUE} \SpecialCharTok{\textasciitilde{}} \ConstantTok{NA\_character\_}
\NormalTok{    ),}
    \AttributeTok{tabagismo\_agrupado =} \FunctionTok{case\_when}\NormalTok{(}
      \CommentTok{\# Grupo de Referência (Protegido/Baixo Risco)}
\NormalTok{      tabagismo }\SpecialCharTok{==} \StringTok{"nunca"} \SpecialCharTok{\textasciitilde{}} \StringTok{"Nunca Fumou"}\NormalTok{,}
      
      \CommentTok{\# Grupo de Risco Passado}
\NormalTok{      tabagismo }\SpecialCharTok{==} \StringTok{"ex{-}consumidor"} \SpecialCharTok{\textasciitilde{}} \StringTok{"Ex{-}Consumidor"}\NormalTok{,}
      
      \CommentTok{\# Grupo de Risco Ativo}
\NormalTok{      tabagismo }\SpecialCharTok{==} \StringTok{"sim"} \SpecialCharTok{\textasciitilde{}} \StringTok{"Consumidor Ativo"}\NormalTok{,}
      
      \CommentTok{\# Tratamento de Missing (s/info, nao avaliado, nti)}
      \ConstantTok{TRUE} \SpecialCharTok{\textasciitilde{}} \ConstantTok{NA\_character\_}
\NormalTok{    )}
\NormalTok{  ) }\SpecialCharTok{\%\textgreater{}\%} 
\NormalTok{  dplyr}\SpecialCharTok{::}\FunctionTok{select}\NormalTok{(}\SpecialCharTok{{-}}\NormalTok{basediag,}\SpecialCharTok{{-}}\NormalTok{estadocivil,}\SpecialCharTok{{-}}\NormalTok{estadofinal,}\SpecialCharTok{{-}}\NormalTok{examesdiagn,}\SpecialCharTok{{-}}\NormalTok{escolaridade,}\SpecialCharTok{{-}}\NormalTok{primeirotrathosp,}\SpecialCharTok{{-}}\NormalTok{ocupacao\_temp,}\SpecialCharTok{{-}}\NormalTok{descricao,}\SpecialCharTok{{-}}\NormalTok{faixa\_idade,}
         \SpecialCharTok{{-}}\NormalTok{origemencaminhamento,}\SpecialCharTok{{-}}\NormalTok{racacor2,}\SpecialCharTok{{-}}\NormalTok{rzntr2,}\SpecialCharTok{{-}}\NormalTok{alcoolismo,}\SpecialCharTok{{-}}\NormalTok{tabagismo,}\SpecialCharTok{{-}}\NormalTok{idade1consulta) }\SpecialCharTok{\%\textgreater{}\%} 
  \FunctionTok{mutate\_if}\NormalTok{(is.character,as.factor)}
\end{Highlighting}
\end{Shaded}

\subsection{Tratamento e Processamento dos
Dados}\label{tratamento-e-processamento-dos-dados}

A base de dados original foi submetida a um rigoroso processo de
limpeza, normalização e \emph{feature engineering} para viabilizar a
análise estatística. O processamento dos dados foi realizado utilizando
a linguagem R, seguindo as etapas descritas abaixo.

\subsubsection{1. Critérios de Inclusão e Limpeza
Inicial}\label{crituxe9rios-de-inclusuxe3o-e-limpeza-inicial}

Foram removidos registros duplicados ou inconsistentes identificados por
artefatos de cabeçalho presentes no corpo dos dados (ex: linhas onde a
variável \texttt{tratamentoantesdodiag} continha o próprio nome da
variável). Registros contendo datas inválidas ou fora do escopo temporal
do estudo foram revisados.

\subsubsection{2. Construção da Variável Desfecho (Resposta ao
Tratamento)}\label{construuxe7uxe3o-da-variuxe1vel-desfecho-resposta-ao-tratamento}

A variável dependente foi construída a partir do campo
\texttt{estadofinal} (``Estado da doença ao final do primeiro
tratamento''). Devido à natureza categórica original, optou-se pela
dicotomização para permitir a regressão logística binária: *
\textbf{Desfecho Satisfatório (1):} Agrupamento das categorias
``Remissão completa'' e ``Remissão parcial''. * \textbf{Desfecho
Insatisfatório (0):} Agrupamento das categorias ``Doença estável'',
``Progressão da doença'', ``Óbito'' e ``Suporte terapêutico''. *
\emph{Exclusão:} Registros classificados como ``Sem informação'' ou
``Não se aplica'' nesta variável foram excluídos da análise de regressão
(análise de casos completos para o desfecho).

\subsubsection{3. Recodificação de Variáveis
Explicativas}\label{recodificauxe7uxe3o-de-variuxe1veis-explicativas}

Para aumentar o poder estatístico e facilitar a interpretação clínica,
as variáveis preditoras foram reclassificadas:

\paragraph{Variáveis
Sociodemográficas}\label{variuxe1veis-sociodemogruxe1ficas}

\begin{itemize}
\tightlist
\item
  \textbf{Raça/Cor:} A variável \texttt{racacor2}, originalmente
  contendo cinco categorias, foi recodificada priorizando os três grupos
  majoritários para análise de desigualdade: ``Branca'', ``Preta'' e
  ``Parda''. As categorias ``Amarela'' e ``Indígena'', devido à baixa
  frequência amostral (\textless1\%), foram tratadas como \emph{missing}
  ou agrupadas conforme a sensibilidade do modelo, sendo a população
  Preta definida como categoria de referência para evidenciar o
  gradiente racial.
\item
  \textbf{Escolaridade:} As oito categorias originais foram reduzidas a
  três níveis funcionais: ``Baixa Escolaridade'' (Analfabeto/Fundamental
  Incompleto), ``Média Escolaridade'' (Fundamental Completo/Médio) e
  ``Alta Escolaridade'' (Superior Incompleto/Completo).
\item
  \textbf{Situação Conjugal:} A variável \texttt{estadocivil} foi
  binariazada em ``Com Companheiro'' (Casados/União Consensual) e ``Sem
  Companheiro'' (Solteiros, Viúvos, Separados), baseando-se na
  literatura sobre suporte social em oncologia.
\end{itemize}

\paragraph{Variáveis Clínicas}\label{variuxe1veis-cluxednicas}

\begin{itemize}
\tightlist
\item
  \textbf{Estadiamento (TNM):} A variável \texttt{estadiam}, contendo
  códigos detalhados (ex: 1, 1C, 2A, 2B), foi simplificada para uma
  variável binária:

  \begin{itemize}
  \tightlist
  \item
    \textbf{Inicial:} Estádios 0, I e II.
  \item
    \textbf{Avançado:} Estádios III e IV.
  \item
    Códigos ``88'' (Não se aplica) e ``99'' (Ignorado) foram convertidos
    para \emph{NA}.
  \end{itemize}
\item
  \textbf{Tratamento Realizado:} A variável \texttt{primeirotrathosp}
  apresentava alta dispersão. Foi reorganizada em quatro modalidades
  principais: ``Cirurgia'', ``Radioterapia'', ``Hormonioterapia'' e
  ``Quimio/Outros''. A categoria ``Sem Tratamento Inicial'' foi criada
  agrupando-se os pacientes classificados como ``Nenhum'' e ``NTI'' (Não
  Teve Indicação), capturando prováveis casos de vigilância ativa ou
  recusa.
\end{itemize}

\paragraph{Classificação Ocupacional (Mineração de
Texto)}\label{classificauxe7uxe3o-ocupacional-minerauxe7uxe3o-de-texto}

A variável \texttt{descricao} (Ocupação baseada na CBO) apresentava 278
categorias distintas de texto livre (ex: ``Administradores\ldots{}'',
``Ajustadores\ldots{}'', ``Trabalhadores da cultura de gramineas''). Foi
aplicada uma técnica de mineração de texto (\emph{string matching}) para
agrupar estas ocupações em cinco macro-categorias socioeconômicas: 1.
\textbf{Rural/Agropecuária:} Palavras-chave como \emph{lavrador, agro,
rural, pesca}. 2. \textbf{Manual/Operacional:} Palavras-chave como
\emph{pedreiro, servente, motorista, mecânico, soldador}. 3.
\textbf{Serviços/Comércio:} Palavras-chave como \emph{vendedor, vigia,
limpeza, doméstico}. 4. \textbf{Técnico/Administrativo:} Palavras-chave
como \emph{técnico, auxiliar, escritório}. 5.
\textbf{Superior/Dirigente:} Palavras-chave como \emph{engenheiro,
médico, professor, diretor}.

\paragraph{Estilo de Vida}\label{estilo-de-vida}

As variáveis \texttt{alcoolismo} e \texttt{tabagismo} mantiveram sua
estrutura tricotômica (``Nunca'', ``Ex-consumidor'', ``Consumidor
Ativo'') para avaliar o efeito do cessaçao e o risco acumulado,
tratando-se ``Sem informação'' como dado ausente.

\subsubsection{\texorpdfstring{4. Tratamento de Dados Ausentes
(\emph{Missing
Data})}{4. Tratamento de Dados Ausentes (Missing Data)}}\label{tratamento-de-dados-ausentes-missing-data}

Foi realizada uma análise de casos completos (\emph{complete-case
analysis}) para a regressão logística. Variáveis com alto percentual de
dados faltantes estruturais, como \texttt{rzntr2} (Razão do não
tratamento), que apresentava \textgreater98\% de preenchimento como
``Não se aplica'', foram excluídas da modelagem multivariada para evitar
instabilidade numérica e viés de convergência.

\subsection{Uma breve visão sobre a distribuição dos dados (após
limpeza)}\label{uma-breve-visuxe3o-sobre-a-distribuiuxe7uxe3o-dos-dados-apuxf3s-limpeza}

\begin{Shaded}
\begin{Highlighting}[]
\FunctionTok{summary}\NormalTok{(dados\_clean)}
\end{Highlighting}
\end{Shaded}

\begin{verbatim}
##           faixadiagtrat         prazoleitratamento      prazoleidiagnostico
##  30 a 60 dias    : 4098   ate 60 dias    : 6564    até 30 dias    :22754   
##  mais de 60 dias :17477   mais de 60 dias:17477    mais de 30 dias: 1287   
##  menos de 30 dias: 2466                                                    
##                                                                            
##                                                                            
##                                                                            
##                                                                            
##  historicofamiliar   estadiam_bin   faixa_etaria_prim_cons origemenSUS
##  nao   :8852       Avançado: 3995   60-69  :9464           Não: 3809  
##  s/info:7872       Inicial :12952   70-79  :8245           Sim:20232  
##  sim   :7317       NA's    : 7094   50-59  :3612                      
##                                     80+    :2265                      
##                                     40-49  : 424                      
##                                     30-39  :  19                      
##                                     (Other):  12                      
##      raca                        basediag_agrupada      status_conjugal_bin
##  Branca: 8585   Confirmacao Microscopica  :23607   Com Companheiro:16247   
##  Parda :10841   Diagnostico Clinico/Imagem:  417   Sem Companheiro: 7098   
##  Preta : 2911   NA's                      :   17   NA's           :  696   
##  NA's  : 1704                                                              
##                                                                            
##                                                                            
##                                                                            
##      resposta_tratamento          tipo_exame_agrupado
##  Insatisfatória:6931     Anatomia Patologica: 9705   
##  Satisfatória  :8224     Marcadores Tumorais:13855   
##  NA's          :8886     NA's               :  481   
##                                                      
##                                                      
##                                                      
##                                                      
##         escolaridade_agrupada             tratamento_agrupado
##  Alta Escolaridade :  803     Cirurgia              :9420    
##  Baixa Escolaridade:13795     Hormonioterapia       :3196    
##  Media Escolaridade: 5498     Quimio/Outros         : 483    
##  NA's              : 3945     Radioterapia          :3968    
##                               Sem Tratamento Inicial:6957    
##                               NA's                  :  17    
##                                                              
##                 classe_trabalho           motivo_nao_tratamento
##  Manual/Operacional     :9263   Condicao Clinica Ruim:   55    
##  Outros/Nao Classificado:1567   Recusa/Abandono      :   24    
##  Rural/Agropecuaria     :6072   Transferido          :   29    
##  Servicos/Comercio      :2006   Tratado/Nao se aplica:23933    
##  Superior/Dirigente     : 615                                  
##  Tecnico/Administrativo : 466                                  
##  NA's                   :4052                                  
##        alcoolismo_agrupado        tabagismo_agrupado
##  Consumidor Ativo:5444     Consumidor Ativo: 3858   
##  Ex-Consumidor   :4484     Ex-Consumidor   : 6365   
##  Nunca Bebeu     :9989     Nunca Fumou     :10078   
##  NA's            :4124     NA's            : 3740   
##                                                     
##                                                     
## 
\end{verbatim}

A análise descritiva da base de dados processada (n = 24.041) revela um
perfil demográfico, clínico e assistencial com características
marcantes.

\subsubsection{Perfil
Sociodemográfico}\label{perfil-sociodemogruxe1fico}

A população estudada é majoritariamente idosa, com concentração
expressiva nas faixas etárias de \textbf{60 a 69 anos} (n=9.464) e
\textbf{70 a 79 anos} (n=8.245), refletindo a história natural da
neoplasia em questão. Em relação à raça/cor, observou-se predominância
de indivíduos autodeclarados \textbf{pardos} (n=10.841), seguidos por
brancos (n=8.585) e pretos (n=2.911).

O perfil socioeconômico aponta para uma população vulnerável: a maioria
absoluta dos pacientes possui \textbf{baixa escolaridade} (n=13.795),
definida como ensino fundamental incompleto ou analfabetismo. A inserção
no mercado de trabalho reflete essa escolaridade, com prevalência de
ocupações \textbf{manuais/operacionais} (n=9.263) e
\textbf{rurais/agropecuárias} (n=6.072). O suporte social parece
robusto, visto que a maior parte dos pacientes vive \textbf{com
companheiro} (n=16.247). A dependência do sistema público é evidente,
com \textbf{84\%} dos pacientes (n=20.232) encaminhados via Sistema
Único de Saúde (SUS).

\subsubsection{Perfil Clínico e
Diagnóstico}\label{perfil-cluxednico-e-diagnuxf3stico}

A confirmação diagnóstica foi predominantemente \textbf{microscópica}
(n=23.607), garantindo alta confiabilidade aos dados. Curiosamente, o
exame mais relevante para o planejamento terapêutico dividiu-se entre
\textbf{Marcadores Tumorais} (n=13.855) e Anatomia Patológica (n=9.705),
sugerindo fortemente o perfil de câncer de próstata, onde o PSA
desempenha papel central na conduta.

Quanto ao estadiamento clínico (TNM), entre os pacientes com informação
válida, houve predominância do \textbf{Estádio Inicial} (n=12.952) em
detrimento do Avançado (n=3.995). Contudo, destaca-se um volume
considerável de dados ausentes (\emph{missing values}) nesta variável
(n=7.094).

\subsubsection{Tratamento e Acesso}\label{tratamento-e-acesso}

Observou-se um gargalo importante no tempo para início do tratamento. A
maioria dos pacientes (n=17.477) iniciou a terapia \textbf{mais de 60
dias} após o diagnóstico, evidenciando desafios no cumprimento da Lei
dos 60 Dias. Por outro lado, o acesso ao diagnóstico foi mais ágil, com
a maioria (n=22.754) sendo diagnosticada em até 30 dias.

A modalidade terapêutica principal foi a \textbf{Cirurgia} (n=9.420),
seguida pela Radioterapia (n=3.968) e Hormonioterapia (n=3.196). O grupo
classificado como ``Sem Tratamento Inicial'' (n=6.957) é expressivo,
podendo representar casos de vigilância ativa (comuns em tumores de
próstata de baixo risco) ou perdas de seguimento.

\subsubsection{Estilo de Vida e
Desfechos}\label{estilo-de-vida-e-desfechos}

O perfil de hábitos de vida mostra que a maioria dos pacientes declarou
``Nunca ter fumado'' (n=10.078) ou ``Nunca ter bebido'' (n=9.989),
embora os grupos de ex-consumidores sejam representativos (Tabagismo:
6.365; Álcool: 4.484).

Em relação ao desfecho primário do estudo, a resposta ao tratamento
(avaliada ao final do primeiro tratamento no hospital) foi
\textbf{Satisfatória} (Remissão completa/parcial) em 8.224 casos e
\textbf{Insatisfatória} (Doença estável/progressão/óbito) em 6.931
casos. Nota-se, entretanto, um alto número de dados não informados para
esta variável (n=8.886), o que justifica a redução amostral observada na
modelagem multivariada subsequente.

\subsection{Regressão Logística}\label{regressuxe3o-loguxedstica}

Para identificar os fatores associados ao desfecho do tratamento,
utilizou-se o modelo de \textbf{Regressão Logística Binomial}. Esta
técnica é o padrão-ouro em estudos epidemiológicos quando a variável
resposta é dicotômica (Sucesso vs.~Falha), pois permite estimar a
probabilidade de ocorrência de um evento (\(P(Y=1)\)) em função de um
conjunto de variáveis explicativas independentes
(\(X_1, X_2, ..., X_k\)).

Matematicamente, o modelo utiliza a função de ligação \emph{logit}, que
transforma a probabilidade em \emph{log-odds} (logaritmo da chance),
linearizando a relação entre as variáveis preditoras e o desfecho:

\[\ln\left(\frac{P}{1-P}\right) = \beta_0 + \beta_1X_1 + \beta_2X_2 + ... + \beta_kX_k\]
Os coeficientes (\(\beta\)) obtidos pelo método da Máxima
Verossimilhança foram exponenciados (\(e^\beta\)) para serem
apresentados como \textbf{Odds Ratios (OR)}, ou Razões de Chances. Um OR
\textgreater{} 1 indica fator de proteção/sucesso, enquanto OR
\textless{} 1 indica fator de risco/falha, considerando um intervalo de
confiança de 95\% (IC95\%).

\subsection{Seleção de Variáveis e Justificativa de
Exclusão}\label{seleuxe7uxe3o-de-variuxe1veis-e-justificativa-de-exclusuxe3o}

A seleção das variáveis para o modelo final seguiu critérios teóricos
(plausibilidade biológica e social) e estatísticos. Inicialmente,
considerou-se a inclusão da variável \textbf{idade} (seja contínua ou em
faixas etárias). No entanto, optou-se por sua exclusão no modelo
multivariado final pelas seguintes razões:

\begin{enumerate}
\def\labelenumi{\arabic{enumi}.}
\tightlist
\item
  \textbf{Homogeneidade da Amostra:} Conforme observado na análise
  descritiva, a distribuição etária da coorte é altamente concentrada
  (mais de 73\% dos pacientes têm entre 60 e 79 anos). Essa baixa
  variabilidade reduz o poder discriminatório da variável idade como
  preditor independente de sucesso terapêutico nesta amostra específica.
\item
  \textbf{Interação com Fatores Sociais:} O foco do estudo recai sobre
  os Determinantes Sociais de Saúde (Raça, Trabalho, Acesso). A idade
  apresenta forte colinearidade com a variável \textbf{Ocupação}
  (pacientes mais idosos tendem a estar na categoria ``Outros/Inativos''
  ou ``Aposentados'', enquanto os mais jovens estão na força de trabalho
  ativa). Manter ambas poderia gerar ruído estatístico e inflacionar
  erros padrão.
\item
  \textbf{História Natural da Doença:} Sendo o câncer (presumivelmente
  de próstata, dado o perfil) uma doença intrinsicamente ligada ao
  envelhecimento, a idade atua como um fator de risco basal universal na
  amostra. Ao removê-la, o modelo purifica o efeito das variáveis de
  interesse (como o racismo estrutural ou barreiras de acesso ao SUS),
  evitando que o peso biológico da idade mascare as desigualdades
  sociais que o estudo visa iluminar.
\end{enumerate}

Dessa forma, o modelo final ajustado priorizou as dimensões de
vulnerabilidade social, estadiamento clínico e modalidades terapêuticas.

\begin{Shaded}
\begin{Highlighting}[]
\NormalTok{dados\_modelo }\OtherTok{\textless{}{-}}\NormalTok{ dados\_clean }\SpecialCharTok{\%\textgreater{}\%}
  \CommentTok{\# 1. Remover quem não tem a informação do desfecho (resposta)}
  \FunctionTok{filter}\NormalTok{(}\SpecialCharTok{!}\FunctionTok{is.na}\NormalTok{(resposta\_tratamento)) }\SpecialCharTok{\%\textgreater{}\%}
  
  \FunctionTok{mutate}\NormalTok{(}
    \CommentTok{\# 2. Criar a variável binária (Y)}
    \CommentTok{\# 1 = Sucesso (Remissão), 0 = Falha (Óbito/Estável/Piora)}
    \AttributeTok{desfecho\_binario =} \FunctionTok{ifelse}\NormalTok{(resposta\_tratamento }\SpecialCharTok{==} \StringTok{"Satisfatória"}\NormalTok{, }\DecValTok{1}\NormalTok{, }\DecValTok{0}\NormalTok{),}
    
    \CommentTok{\# 3. Ajustar as Referências (Baselines)}
    \CommentTok{\# Isso define quem é o "1.0" na comparação de risco.}
    
    \CommentTok{\# Sexo/Raça (Geralmente a maioria ou o grupo de risco padrão)}
    \AttributeTok{raca =} \FunctionTok{relevel}\NormalTok{(}\FunctionTok{as.factor}\NormalTok{(raca), }\AttributeTok{ref =} \StringTok{"Preta"}\NormalTok{), }
    
    \CommentTok{\# Idade (Se estiver usando contínua, ok. Se for faixa, defina a menor)}
    \CommentTok{\# (Assumindo que você manteve idade1consulta como numérica)}
    
    \CommentTok{\# Escolaridade: Referência = Baixa}
    \AttributeTok{escolaridade\_agrupada =} \FunctionTok{relevel}\NormalTok{(}\FunctionTok{as.factor}\NormalTok{(escolaridade\_agrupada), }\AttributeTok{ref =} \StringTok{"Alta Escolaridade"}\NormalTok{),}
    
    \CommentTok{\# Estádio: Referência = Inicial}
    \AttributeTok{estadiam\_bin =} \FunctionTok{relevel}\NormalTok{(}\FunctionTok{as.factor}\NormalTok{(estadiam\_bin), }\AttributeTok{ref =} \StringTok{"Avançado"}\NormalTok{),}
    
    \CommentTok{\# Tratamento: Referência = Sem Tratamento (Para ver o quanto a Cirurgia protege)}
    \CommentTok{\# OU Referência = Cirurgia (Para ver o quanto não tratar piora) {-}\textgreater{} Vou sugerir Cirurgia}
    \AttributeTok{tratamento\_agrupado =} \FunctionTok{relevel}\NormalTok{(}\FunctionTok{as.factor}\NormalTok{(tratamento\_agrupado), }\AttributeTok{ref =} \StringTok{"Sem Tratamento Inicial"}\NormalTok{),}
    
    \CommentTok{\# Vícios: Referência = Nunca}
    \AttributeTok{alcoolismo\_agrupado =} \FunctionTok{relevel}\NormalTok{(}\FunctionTok{as.factor}\NormalTok{(alcoolismo\_agrupado), }\AttributeTok{ref =} \StringTok{"Consumidor Ativo"}\NormalTok{),}
    \AttributeTok{tabagismo\_agrupado =} \FunctionTok{relevel}\NormalTok{(}\FunctionTok{as.factor}\NormalTok{(tabagismo\_agrupado), }\AttributeTok{ref =} \StringTok{"Consumidor Ativo"}\NormalTok{),}
    
    \CommentTok{\# Classe Trabalho: Referência = Manual ou Rural}
    \AttributeTok{classe\_trabalho =} \FunctionTok{relevel}\NormalTok{(}\FunctionTok{as.factor}\NormalTok{(classe\_trabalho), }\AttributeTok{ref =} \StringTok{"Outros/Nao Classificado"}\NormalTok{),}
    
    \CommentTok{\# Classe Trabalho: Referência = Manual ou Rural}
    \AttributeTok{origemenSUS =} \FunctionTok{relevel}\NormalTok{(}\FunctionTok{as.factor}\NormalTok{(origemenSUS), }\AttributeTok{ref =} \StringTok{"Sim"}\NormalTok{)}
\NormalTok{  )}

\DocumentationTok{\#\# MODELO }

\NormalTok{modelo\_logistico }\OtherTok{\textless{}{-}} \FunctionTok{glm}\NormalTok{(}
\NormalTok{  desfecho\_binario }\SpecialCharTok{\textasciitilde{}} 
    \CommentTok{\#faixa\_etaria\_prim\_cons +   \# Faixa Etária na Primeira Consulta}
\NormalTok{    raca }\SpecialCharTok{+}                \CommentTok{\# Raça}
\NormalTok{    status\_conjugal\_bin }\SpecialCharTok{+}      \CommentTok{\# Suporte social}
\NormalTok{    escolaridade\_agrupada }\SpecialCharTok{+}    \CommentTok{\# Nível socioeconômico}
\NormalTok{    classe\_trabalho }\SpecialCharTok{+}          \CommentTok{\# Exposição ocupacional}
\NormalTok{    estadiam\_bin }\SpecialCharTok{+}             \CommentTok{\# Gravidade da doença (Usei a binária Inicial/Avançado)}
\NormalTok{    tratamento\_agrupado }\SpecialCharTok{+}      \CommentTok{\# O que foi feito}
\NormalTok{    alcoolismo\_agrupado }\SpecialCharTok{+}      \CommentTok{\# Comorbidade 1}
\NormalTok{    tabagismo\_agrupado }\SpecialCharTok{+}       \CommentTok{\# Comorbidade 2}
\NormalTok{    origemenSUS,               }\CommentTok{\# Origem do paciente}
  
  \AttributeTok{family =} \FunctionTok{binomial}\NormalTok{(}\AttributeTok{link =} \StringTok{"logit"}\NormalTok{),}
  \AttributeTok{data =}\NormalTok{ dados\_modelo}
\NormalTok{)}

\CommentTok{\# Resumo estatístico bruto}
\FunctionTok{summary}\NormalTok{(modelo\_logistico) }\CommentTok{\# O MODELO COMPLETO É TB O MELHOR MODELO. TODOS OS VIFs DERAM MENOR QUE 5}
\end{Highlighting}
\end{Shaded}

\begin{verbatim}
## 
## Call:
## glm(formula = desfecho_binario ~ raca + status_conjugal_bin + 
##     escolaridade_agrupada + classe_trabalho + estadiam_bin + 
##     tratamento_agrupado + alcoolismo_agrupado + tabagismo_agrupado + 
##     origemenSUS, family = binomial(link = "logit"), data = dados_modelo)
## 
## Coefficients:
##                                         Estimate Std. Error z value Pr(>|z|)
## (Intercept)                             -0.27381    0.20153  -1.359 0.174260
## racaBranca                              -0.24083    0.07413  -3.249 0.001159
## racaParda                                0.38053    0.07585   5.017 5.26e-07
## status_conjugal_binSem Companheiro       0.09173    0.05418   1.693 0.090477
## escolaridade_agrupadaBaixa Escolaridade  0.26284    0.16205   1.622 0.104802
## escolaridade_agrupadaMedia Escolaridade  0.30874    0.16218   1.904 0.056948
## classe_trabalhoManual/Operacional       -0.66942    0.09095  -7.360 1.84e-13
## classe_trabalhoRural/Agropecuaria       -0.24702    0.09267  -2.665 0.007689
## classe_trabalhoServicos/Comercio        -0.30203    0.11554  -2.614 0.008949
## classe_trabalhoSuperior/Dirigente       -0.42356    0.19488  -2.173 0.029745
## classe_trabalhoTecnico/Administrativo   -0.50125    0.18857  -2.658 0.007858
## estadiam_binInicial                      0.77568    0.06027  12.870  < 2e-16
## tratamento_agrupadoCirurgia              1.02710    0.07407  13.867  < 2e-16
## tratamento_agrupadoHormonioterapia      -0.81283    0.07042 -11.543  < 2e-16
## tratamento_agrupadoQuimio/Outros        -1.64570    0.25018  -6.578 4.77e-11
## tratamento_agrupadoRadioterapia         -0.27178    0.06989  -3.889 0.000101
## alcoolismo_agrupadoEx-Consumidor         0.16660    0.07707   2.162 0.030647
## alcoolismo_agrupadoNunca Bebeu          -0.16026    0.06413  -2.499 0.012452
## tabagismo_agrupadoEx-Consumidor          0.35534    0.07533   4.717 2.39e-06
## tabagismo_agrupadoNunca Fumou            0.19209    0.07113   2.701 0.006922
## origemenSUSNão                           0.38259    0.08358   4.577 4.71e-06
##                                            
## (Intercept)                                
## racaBranca                              ** 
## racaParda                               ***
## status_conjugal_binSem Companheiro      .  
## escolaridade_agrupadaBaixa Escolaridade    
## escolaridade_agrupadaMedia Escolaridade .  
## classe_trabalhoManual/Operacional       ***
## classe_trabalhoRural/Agropecuaria       ** 
## classe_trabalhoServicos/Comercio        ** 
## classe_trabalhoSuperior/Dirigente       *  
## classe_trabalhoTecnico/Administrativo   ** 
## estadiam_binInicial                     ***
## tratamento_agrupadoCirurgia             ***
## tratamento_agrupadoHormonioterapia      ***
## tratamento_agrupadoQuimio/Outros        ***
## tratamento_agrupadoRadioterapia         ***
## alcoolismo_agrupadoEx-Consumidor        *  
## alcoolismo_agrupadoNunca Bebeu          *  
## tabagismo_agrupadoEx-Consumidor         ***
## tabagismo_agrupadoNunca Fumou           ** 
## origemenSUSNão                          ***
## ---
## Signif. codes:  0 '***' 0.001 '**' 0.01 '*' 0.05 '.' 0.1 ' ' 1
## 
## (Dispersion parameter for binomial family taken to be 1)
## 
##     Null deviance: 10453.8  on 7788  degrees of freedom
## Residual deviance:  9282.2  on 7768  degrees of freedom
##   (7366 observations deleted due to missingness)
## AIC: 9324.2
## 
## Number of Fisher Scoring iterations: 4
\end{verbatim}

\subsection{Os resultados do modelo}\label{os-resultados-do-modelo}

A análise multivariada por regressão logística foi conduzida para
identificar os determinantes da resposta satisfatória ao tratamento
(tabela abaixo). O modelo final demonstrou-se robusto, identificando
associações significativas em dimensões sociodemográficas, clínicas e
comportamentais.

\begin{Shaded}
\begin{Highlighting}[]
\CommentTok{\# Gerar tabela}
\NormalTok{tabela\_publicacao }\OtherTok{\textless{}{-}} \FunctionTok{tbl\_regression}\NormalTok{(modelo\_logistico, }
                                    \AttributeTok{exponentiate =} \ConstantTok{TRUE}\NormalTok{, }\CommentTok{\# CRUCIAL: Converte para Odds Ratio}
                                    \AttributeTok{pvalue\_fun =} \SpecialCharTok{\textasciitilde{}}\FunctionTok{style\_pvalue}\NormalTok{(.x, }\AttributeTok{digits =} \DecValTok{3}\NormalTok{),}
                                    \AttributeTok{label =} \FunctionTok{list}\NormalTok{(}
                                      \CommentTok{\#faixa\_etaria\_prim\_cons \textasciitilde{} "Faixa Etária (anos)",}
\NormalTok{                                      raca }\SpecialCharTok{\textasciitilde{}} \StringTok{"Raça"}\NormalTok{,}
\NormalTok{                                      status\_conjugal\_bin }\SpecialCharTok{\textasciitilde{}} \StringTok{"Situação Conjugal"}\NormalTok{,}
\NormalTok{                                      escolaridade\_agrupada }\SpecialCharTok{\textasciitilde{}} \StringTok{"Escolaridade"}\NormalTok{,}
\NormalTok{                                      classe\_trabalho }\SpecialCharTok{\textasciitilde{}} \StringTok{"Ocupação"}\NormalTok{,}
\NormalTok{                                      estadiam\_bin }\SpecialCharTok{\textasciitilde{}} \StringTok{"Estadiamento"}\NormalTok{,}
\NormalTok{                                      tratamento\_agrupado }\SpecialCharTok{\textasciitilde{}} \StringTok{"Tratamento Realizado"}\NormalTok{,}
\NormalTok{                                      alcoolismo\_agrupado }\SpecialCharTok{\textasciitilde{}} \StringTok{"Histórico de Álcool"}\NormalTok{,}
\NormalTok{                                      tabagismo\_agrupado }\SpecialCharTok{\textasciitilde{}} \StringTok{"Histórico de Fumo"}\NormalTok{,}
\NormalTok{                                      origemenSUS }\SpecialCharTok{\textasciitilde{}} \StringTok{"Encaminhado pelo SUS"}
\NormalTok{                                    )) }\SpecialCharTok{\%\textgreater{}\%}
  \FunctionTok{bold\_p}\NormalTok{(}\AttributeTok{t =} \FloatTok{0.05}\NormalTok{) }\SpecialCharTok{\%\textgreater{}\%} \CommentTok{\# Negrita quem for significativo}
  \FunctionTok{bold\_labels}\NormalTok{() }\SpecialCharTok{\%\textgreater{}\%}    \CommentTok{\# Negrita os nomes das variáveis}
  \FunctionTok{italicize\_levels}\NormalTok{()   }\CommentTok{\# Coloca as categorias em itálico}

\CommentTok{\# Visualizar}
\NormalTok{tabela\_publicacao}
\end{Highlighting}
\end{Shaded}

\begin{table}[t]
\fontsize{12.0pt}{14.0pt}\selectfont
\begin{tabular*}{\linewidth}{@{\extracolsep{\fill}}lccc}
\toprule
\textbf{Characteristic} & \textbf{OR} & \textbf{95\% CI} & \textbf{p-value} \\ 
\midrule\addlinespace[2.5pt]
{\bfseries Ra\c{c}a} &  &  &  \\ 
{\itshape     Preta} & \textemdash & \textemdash &  \\ 
{\itshape     Branca} & 0.79 & 0.68, 0.91 & {\bfseries 0.001} \\ 
{\itshape     Parda} & 1.46 & 1.26, 1.70 & {\bfseries <0.001} \\ 
{\bfseries Situa\c{c}\~ao Conjugal} &  &  &  \\ 
{\itshape     Com Companheiro} & \textemdash & \textemdash &  \\ 
{\itshape     Sem Companheiro} & 1.10 & 0.99, 1.22 & 0.090 \\ 
{\bfseries Escolaridade} &  &  &  \\ 
{\itshape     Alta Escolaridade} & \textemdash & \textemdash &  \\ 
{\itshape     Baixa Escolaridade} & 1.30 & 0.95, 1.79 & 0.105 \\ 
{\itshape     Media Escolaridade} & 1.36 & 0.99, 1.87 & 0.057 \\ 
{\bfseries Ocupa\c{c}\~ao} &  &  &  \\ 
{\itshape     Outros/Nao Classificado} & \textemdash & \textemdash &  \\ 
{\itshape     Manual/Operacional} & 0.51 & 0.43, 0.61 & {\bfseries <0.001} \\ 
{\itshape     Rural/Agropecuaria} & 0.78 & 0.65, 0.94 & {\bfseries 0.008} \\ 
{\itshape     Servicos/Comercio} & 0.74 & 0.59, 0.93 & {\bfseries 0.009} \\ 
{\itshape     Superior/Dirigente} & 0.65 & 0.45, 0.96 & {\bfseries 0.030} \\ 
{\itshape     Tecnico/Administrativo} & 0.61 & 0.42, 0.88 & {\bfseries 0.008} \\ 
{\bfseries Estadiamento} &  &  &  \\ 
{\itshape     Avan\c{c}ado} & \textemdash & \textemdash &  \\ 
{\itshape     Inicial} & 2.17 & 1.93, 2.44 & {\bfseries <0.001} \\ 
{\bfseries Tratamento Realizado} &  &  &  \\ 
{\itshape     Sem Tratamento Inicial} & \textemdash & \textemdash &  \\ 
{\itshape     Cirurgia} & 2.79 & 2.42, 3.23 & {\bfseries <0.001} \\ 
{\itshape     Hormonioterapia} & 0.44 & 0.39, 0.51 & {\bfseries <0.001} \\ 
{\itshape     Quimio/Outros} & 0.19 & 0.12, 0.31 & {\bfseries <0.001} \\ 
{\itshape     Radioterapia} & 0.76 & 0.66, 0.87 & {\bfseries <0.001} \\ 
{\bfseries Hist\\'orico de \'Alcool} &  &  &  \\ 
{\itshape     Consumidor Ativo} & \textemdash & \textemdash &  \\ 
{\itshape     Ex-Consumidor} & 1.18 & 1.02, 1.37 & {\bfseries 0.031} \\ 
{\itshape     Nunca Bebeu} & 0.85 & 0.75, 0.97 & {\bfseries 0.012} \\ 
{\bfseries Hist\\'orico de Fumo} &  &  &  \\ 
{\itshape     Consumidor Ativo} & \textemdash & \textemdash &  \\ 
{\itshape     Ex-Consumidor} & 1.43 & 1.23, 1.65 & {\bfseries <0.001} \\ 
{\itshape     Nunca Fumou} & 1.21 & 1.05, 1.39 & {\bfseries 0.007} \\ 
{\bfseries Encaminhado pelo SUS} &  &  &  \\ 
{\itshape     Sim} & \textemdash & \textemdash &  \\ 
{\itshape     N\~ao} & 1.47 & 1.25, 1.73 & {\bfseries <0.001} \\ 
\bottomrule
\end{tabular*}
\begin{minipage}{\linewidth}
Abbreviations: CI = Confidence Interval, OR = Odds Ratio\\
\end{minipage}
\end{table}

No tocante às características raciais, utilizando a população preta como
grupo de referência, observou-se um padrão heterogêneo. Pacientes
autodeclarados pardos apresentaram uma chance \textbf{46\% maior} de
resposta satisfatória (OR: 1,46; IC95\%: 1,26--1,70; p\textless0,001).
Em contrapartida, a população branca apresentou uma chance reduzida de
sucesso terapêutico em comparação à população preta (OR: 0,79; IC95\%:
0,68--0,91; p=0,001).

As desigualdades socioeconômicas manifestaram-se fortemente através da
ocupação e da origem do encaminhamento. Pacientes não encaminhados pelo
Sistema Único de Saúde (SUS) --- atendidos via saúde suplementar ou
particular --- tiveram 1,47 vezes a chance de obter resposta favorável
comparados aos pacientes do SUS (IC95\%: 1,25--1,73; p\textless0,001).
Quanto à ocupação, tendo como referência o grupo ``Outros/Não
Classificado'' (que possivelmente engloba inativos com renda ou
categorias não manuais não listadas), todas as classes laborais
específicas apresentaram prognóstico inferior. Destaca-se a categoria
``Manual/Operacional'' com a menor chance de resposta (OR: 0,51; IC95\%:
0,43--0,61), evidenciando a vulnerabilidade de trabalhadores braçais. A
escolaridade e a situação conjugal não atingiram significância
estatística no nível de 5\%.

As variáveis clínicas comportaram-se conforme o esperado biologicamente.
O diagnóstico em estádio inicial mais que dobrou a chance de sucesso
(OR: 2,17; IC95\%: 1,93--2,44) em relação ao estádio avançado. A escolha
terapêutica foi determinante: comparado aos pacientes sem tratamento
inicial (provável vigilância ativa ou recusa), a cirurgia elevou a
chance de resposta satisfatória em \textbf{2,79 vezes}
(p\textless0,001). Inversamente, modalidades isoladas como
Hormonioterapia (OR: 0,44) e Quimioterapia (OR: 0,19) associaram-se a
desfechos inferiores à referência, refletindo provável doença sistêmica
ou inoperável nestes grupos.

Em relação ao estilo de vida, o histórico de tabagismo mostrou que
ex-fumantes (OR: 1,43) e não fumantes (OR: 1,21) possuem melhor
prognóstico que fumantes ativos. Para o etilismo, observou-se um
aparente paradoxo: enquanto ex-consumidores tiveram melhor resposta que
consumidores ativos (OR: 1,18), o grupo que ``Nunca Bebeu'' apresentou
chance inferior aos que bebem (OR: 0,85; p=0,012).

\subsection{Discussão}\label{discussuxe3o}

Os resultados corroboram a literatura sobre a influência dos
determinantes sociais na saúde (SDSS), mas apresentam particularidades
regionais importantes. A vantagem de sobrevivência observada na
população parda em relação à preta alinha-se a estudos brasileiros que
apontam para o ``gradiente de cor'', onde o racismo estrutural e
barreiras de acesso incidem com maior severidade sobre a população preta
fenotipicamente mais marcada (LOPEZ, 2011). O achado de que brancos
tiveram desempenho inferior aos pretos nesta amostra específica diverge
de dados globais, sugerindo possíveis viéses locais de seleção, onde a
população branca pode estar chegando ao serviço com idade mais avançada
ou comorbidades não mensuradas no modelo.

A disparidade baseada no tipo de encaminhamento (SUS versus Não-SUS)
reitera a fragmentação do sistema de saúde brasileiro. A ``Lei dos 60
Dias'' busca mitigar atrasos, mas pacientes da rede privada
historicamente acessam métodos diagnósticos e terapêuticos com maior
celeridade, impactando diretamente o estadiamento e a sobrevida (SOUZA
et al., 2020). O impacto da ocupação reforça essa tese: trabalhadores
manuais e rurais, frequentemente expostos a riscos ocupacionais e com
menor flexibilidade de horário para consultas, apresentaram os piores
desfechos.

Clinicamente, a superioridade da cirurgia confirma o tratamento
cirúrgico como padrão-ouro para doença localizada com intenção curativa.
O baixo desempenho das terapias sistêmicas (hormônio e quimioterapia)
reflete o viés de indicação, onde tais modalidades são reservadas para
casos onde a cura cirúrgica já não é viável (tumor avançado ou
metastático).

Por fim, o achado paradoxal referente ao consumo de álcool --- onde
abstêmios tiveram pior prognóstico que bebedores ativos --- pode ser
explicado pelo ``efeito do abstêmio doente'' (\emph{sick quitter
effect}). Este viés epidemiológico sugere que o grupo de ``nunca
bebedores'' frequentemente inclui indivíduos com saúde frágil ou
comorbidades prévias que os impediram de iniciar o consumo de álcool,
enquanto bebedores ativos moderados tendem a possuir vida social ativa e
melhor estado funcional geral (SHAPER et al., 1988; STOCKWELL et al.,
2016). Já o tabagismo mantém sua coerência biológica nociva, onde a
cessação ou não exposição confirma-se como fator protetor robusto.

\subsection{Referências}\label{referuxeancias}

\phantomsection\label{refs}
\begin{CSLReferences}{0}{1}
\end{CSLReferences}

\end{document}
